\section{第9组：气体分子运动论}
\sourceproblem{9-T1}{55}{59}
\problemtext 氧气瓶容积 $3.2\times10^{-2}$ 立方米，初压 $1.30\times10^7$ Pa，补气阈值 $1.0\times10^6$ Pa，每日耗用常压 $1.01\times10^5$ Pa 下0.40立方米氧气，等温求可用天数。
\workings 由 $pV=Nk_BT$，可消耗分子数为 $(p_i-p_f)V/(k_BT)$，每日用量为 $p_dV_d/(k_BT)$。故
\[
t=\frac{(1.30\times10^7-1.0\times10^6)0.032}{1.01\times10^5(0.40)}\approx9.50\ \text{天}.
\]
\crosscheck 不可将瓶内剩余气体全部计为可用；高压氧气的非理想性按教材理想模型忽略。
\sourceproblem{9-T2}{55}{59}
\problemtext 真空压强 $1.01325\times10^{-10}$ Pa、温度25摄氏度，求数密度，并与标准状况比较。
\workings $n=p/(k_BT)$，取 $k_B=1.380649\times10^{-23}$ J/K，$T=298.15$ K，得 $n\approx2.46\times10^{10}\,\mathrm{m^{-3}}$。标准状况取101325 Pa和273.15 K，$n_0\approx2.69\times10^{25}\,\mathrm{m^{-3}}$，约为前者的 $1.09\times10^{15}$ 倍。
\crosscheck 已放大原页确认真空压强指数负10；真空并非完全没有分子。
\sourceproblem{9-T3}{56}{60}
\problemtext 解释原题八个关于速率分布函数的式子。
\workings 设总分子数 $N$、速率分布密度 $f(v)$、区间 $A=[v_1,v_2]$，单分子质量 $m$。依原书顺序：(1) $f(v)=\mathrm dN/(N\mathrm dv)$ 为速率概率密度；(2) $f(v)\mathrm dv$ 为微小速率区间的分子比例；(3) $Nf(v)\mathrm dv$ 为该区间分子数；(4) $\int_A f(v)\mathrm dv$ 为有限区间比例；(5) $N\int_A f(v)\mathrm dv$ 为该区间分子数；(6) $\int_A vf(v)\mathrm dv/\int_A f(v)\mathrm dv$ 为区间条件平均速率；(7) $N\int_A mv^2f(v)\mathrm dv/2$ 为区间分子平动动能总和；(8) $\int_A mv^2f(v)\mathrm dv/[2\int_A f(v)\mathrm dv]$ 为该区间每分子平均平动动能。
\crosscheck 条件平均的分母必须非零；$f$ 的量纲为速率的倒数，不能把概率密度本身当概率。
\sourceproblem{9-T4}{56}{60}
\problemtext 同温氢气和氧气的麦克斯韦速率分布，右峰最概然速率2000 m/s，求两气体最概然速率及温度。
\workings $v_p=\sqrt{2RT/M}$，较轻的氢气对应右峰II，氧气对应I。摩尔质量比16，故 $v_{p,\mathrm O_2}=500$ m/s。取氢气摩尔质量 $0.002$ kg/mol，$T=Mv_p^2/(2R)\approx481$ K。
\crosscheck 图中横坐标2000已核实；不由示意图左峰位置直接估读，使用质量平方根关系。
\sourceproblem{9-T5}{57}{61}
\problemtext 气体方均根速率450 m/s、压强 $7\times10^4$ Pa，求质量密度。
\workings 气体动理论 $p=\rho\overline{v^2}/3$，故 $\rho=3p/(450)^2\approx1.037\,\mathrm{kg/m^3}$。
\crosscheck 方均根速率的平方才是均方速率。
\sourceproblem{9-T6}{57}{61}
\problemtext 氧气压强 $1.01\times10^5$ Pa、温度27摄氏度，求分子数密度、质量密度、分子平均平动动能。
\workings $T=300.15$ K，$n=p/(k_BT)\approx2.44\times10^{25}\,\mathrm{m^{-3}}$；取 $M=0.032$ kg/mol，$\rho=pM/(RT)\approx1.30\,\mathrm{kg/m^3}$；$\overline K=3k_BT/2\approx6.22\times10^{-21}$ J。
\crosscheck 平动能只含三个平动自由度，不使用双原子总能量的5/2系数。
\sourceproblem{9-T7}{57}{61}
\problemtext 容积 $1.0\times10^{-3}$ 立方米内有 $1.01\times10^{23}$ 个氢气分子，压强 $1.01\times10^5$ Pa，求温度和方均根速率。
\workings $T=pV/(Nk_B)\approx72.43$ K，$v_{\rm rms}=\sqrt{3RT/M}\approx950$ m/s（$M=0.002$ kg/mol）。
\crosscheck 原页三个指数已高清核实；可直接由 $v_{\rm rms}^2=3pV/(Nm)$ 独立验证。
\sourceproblem{9-T8}{58}{62}
\problemtext 刚性双原子理想气体体积 $2.0\times10^{-3}$ 立方米、内能675 J，分子数 $5.4\times10^{22}$，求压强、平均平动动能和温度。
\workings 刚性双原子 $U=5Nk_BT/2=5pV/2$，故 $p=1.35\times10^5$ Pa；每分子平均平动能 $3U/(5N)=7.50\times10^{-21}$ J，$T=2U/(5Nk_B)\approx362$ K。
\crosscheck 按题给刚性模型忽略振动；内能与平动能不是同一概念。
\sourceproblem{9-T9}{58}{62}
\problemtext 1 mol双原子气体在27摄氏度，求全体分子的平动与转动动能总和。
\workings 刚性双原子分子有三个平动、两个转动自由度。$E_{\rm tr}=3RT/2\approx3.74\times10^3$ J，$E_{\rm rot}=RT\approx2.50\times10^3$ J。
\crosscheck 原题问1 mol的总贡献，不是每分子能量；若求每分子则将 $R$ 改为 $k_B$。
\sourceproblem{9-T10}{58}{62}
\problemtext 1 mol水蒸气分解为同温1 mol氢气和半mol氧气，不计振动自由度，求内能增长百分比。
\workings 非线性水分子有6个平动转动自由度，热运动内能 $U_i=3RT$。产物为1.5 mol刚性双原子气体，$U_f=1.5(5RT/2)=15RT/4$，故 $(U_f-U_i)/U_i=25\%$。
\crosscheck 25\%仅为教材自由度模型的热运动内能变化；实际化学分解涉及键能和化学内能，不能用该百分比代替反应内能。

\sourceproblem{9-T11}{59}{63}
\problemtext 刚性分子自由度 $i=t+r$，气体质量 $m$、摩尔质量 $M$，解释原题六个能量表达式。
\workings 按原序：(1) $k_BT/2$ 是单个二次型自由度的平均能量；(2) $3k_BT/2$ 是每分子平均平动动能；(3) $ik_BT/2$ 是每分子平均平动及转动动能；(4) $iRT/2$ 是每摩尔热运动内能；(5) $(m/M)3RT/2$ 是全体分子平动动能；(6) $(m/M)iRT/2$ 是该气体热运动内能。
\crosscheck 采用经典能量均分、理想气体及无振动的刚性模型；$m/M$ 的量纲为摩尔。
\sourceproblem{9-T12}{59}{63}
\problemtext $N$ 个质量均为 $m$ 的分子，图示数密度函数 $Nf(v)$ 从0到 $v_0$ 线性升至 $a$，再至 $2v_0$ 保持 $a$，其余为零。求面积意义、$a$、区间 $[v_0/2,3v_0/2]$ 内分子数及平均平动动能。
\workings 图下面积为总分子数，$N=av_0/2+av_0=3av_0/2$，故 $a=2N/(3v_0)$。指定区间分子数
\[
N_A=\int_{v_0/2}^{v_0}\frac{av}{v_0}\mathrm dv+\int_{v_0}^{3v_0/2}a\mathrm dv
=\frac78av_0=\frac7{12}N.
\]
平均平动动能
\[
\overline K=\frac{ma}{2N}\left[\frac1{v_0}\int_0^{v_0}v^3\mathrm dv+\int_{v_0}^{2v_0}v^2\mathrm dv\right]
=\frac{31}{36}mv_0^2.
\]
\crosscheck 归一化积分为 $N$，指定区间比例小于1，$\overline{v^2}=31v_0^2/18$ 小于最大速率平方 $4v_0^2$。
